% Generated draft: requires mathematical and source review.
\documentclass[11pt]{article}
\usepackage[margin=25mm]{geometry}
\usepackage{fontspec}
\setmainfont{Latin Modern Roman}
\usepackage{amsmath,amssymb,mathtools,unicode-math}
\setmathfont{Latin Modern Math}
\usepackage{enumitem,needspace,xurl,hyperref}
\setlength{\parindent}{0pt}
\setlength{\parskip}{5pt}
\setlength{\emergencystretch}{3em}
\newcommand{\N}{\mathbb{N}}
\newcommand{\Z}{\mathbb{Z}}
\newcommand{\Q}{\mathbb{Q}}
\newcommand{\R}{\mathbb{R}}
\newcommand{\C}{\mathbb{C}}
\newcommand{\F}{\mathbb{F}}
\newcommand{\A}{\mathbb{A}}
\newcommand{\PP}{\mathbb{P}}
\newcommand{\E}{\mathbb{E}}
\newcommand{\eps}{\varepsilon}
\newcommand{\dd}{\,\mathrm{d}}
\begin{document}
\section*{UnsolvedMath ID: 30004365}
\textbf{Source problem code:} catalogue entry
\subsection{Definitions and mathematical statement}
Let \(k\) denote the ground field required by the source theorem. For a quiver \(Q\), its path algebra \(kQ\) has directed paths as a basis, with multiplication given by concatenation when defined and zero otherwise. An ideal of relations \(I\subseteq kQ\) gives the presented algebra \(kQ/I\). Write
\[
A(\Delta)=kQ(\Delta)/I(\Delta).
\]
The input is the quiver with relations \((Q(\Delta),I(\Delta))\), promised to define a gentle algebra in the class covered by the source's Theorem 3.3. The notation \(\Delta\) does not supply additional geometric input.

In the usual finite-dimensional setting, derived equivalence means equivalence of the bounded derived categories of finite-dimensional modules as \(k\)-linear triangulated categories. The supplied record does not reproduce the ground-field assumptions, the precise definition and scope of gentle algebras used in Theorem 3.3, or that theorem's convention for derived equivalence; these must be recovered before fixing the domain precisely.

\textbf{Requested algorithm, conditional on the missing theorem specification.} Let \(\mathcal C_{3.3}\) denote the intended class of gentle algebras and let \(\operatorname{Inv}_{3.3}(A)\) denote the complete derived-equivalence invariant of the referenced Theorem 3.3. These symbols are placeholders for the missing theorem data. Construct an algorithm which, for every valid input presentation \((Q,I)\) with \(kQ/I\in\mathcal C_{3.3}\), terminates and computes \(\operatorname{Inv}_{3.3}(kQ/I)\) using only that presentation. Its output must include all of the theorem's genus, puncture and boundary data, winding-number invariants, and remaining parity or greatest-common-divisor invariants, with their applicable case distinctions and identifications. In particular, after the theorem's output conventions have been specified, completeness requires
\[
\operatorname{Inv}_{3.3}(A)\sim\operatorname{Inv}_{3.3}(B)
\quad\Longleftrightarrow\quad
A\text{ and }B\text{ are derived equivalent}
\qquad(A,B\in\mathcal C_{3.3}),
\]
where \(\sim\) means agreement under the identifications prescribed by that theorem.

\textbf{Unresolved specification.} The supplied record does not reproduce Theorem 3.3 or unambiguously locate it within a particular supplied reference. Consequently, it does not determine the exact output tuple, winding-number conventions, boundary and puncture conventions, parity and greatest-common-divisor formulas, or their case restrictions. The displayed requirement records the intended target; it is not a fully specified computational problem until these omissions are resolved.
\subsection{Short English statement}
For UnsolvedMath ID \(30004365\), construct a terminating procedure that extracts the complete derived-equivalence invariant of a gentle algebra directly from its quiver and relations. It must recover the specified topological, winding-number, parity and greatest-common-divisor data without additional geometric input. The referenced Theorem 3.3 is absent, so its hypotheses and exact output specification remain unresolved. The supplied record labels the problem open; this draft does not independently verify that status.
\subsection{Sources}
\begin{itemize}
\item Supplied local UnsolvedMath registry record: permanent ID \(30004365\), title \emph{Derived Invariants from Gentle Quivers with Relations}, source problem code \texttt{catalogue entry}, collection \emph{Oberwolfach Reports Open Problems}. The record gives proposed year \(2020\) and names no proposer.
\item Supplied report citation: \emph{Representation Theory of Quivers and Finite Dimensional Algebras} (2021), Oberwolfach Reports, pp.\ 143--230, DOI \texttt{10.4171/owr/2020/3}. \url{https://doi.org/10.4171/owr/2020/3}. The supplied background describes these reports as documenting mathematical workshops and their open-problem discussions. The report year and DOI are retained as supplied.
\item Supplied literature source: \emph{A complete derived invariant for gentle algebras}. \url{https://doi.org/10.1007/s00029-022-00822-x}. No author information or theorem text is supplied in the record.
\item Literature assessment dated exactly \texttt{2026-08-22}: the supplied curation labels the problem \emph{open} and states that a complete derived invariant is known abstractly for gentle algebras, while a general effective algorithm extracting every invariant from the quiver with relations is still not supplied. This is supplied dated curation derived from the OWR material, not an independently verified literature assessment. Its time-sensitive status and cited claims require independent verification.
\item The registry labels the statement \texttt{corrected\_verified} and reports a manual semantic, typographic and LaTeX review asserting that the clean statement is self-contained and preserves the target. This is supplied review metadata; it does not resolve the missing specification of Theorem 3.3 in this draft.
\end{itemize}
\end{document}
